\documentclass[11pt, a4paper]{article}

\usepackage{bm}
\usepackage[utf8]{inputenc}
\usepackage[italian]{babel}
\usepackage{amsmath, amssymb}
\usepackage{amsthm}
\usepackage{graphicx}
\usepackage{booktabs}
\usepackage{geometry}
\geometry{a4paper, top=2.5cm, bottom=2.5cm, left=2.5cm, right=2.5cm}
\usepackage{caption}
\captionsetup{justification=centering}
\usepackage{float}
\usepackage[hidelinks]{hyperref}
\usepackage{subcaption}
\usepackage{verbatim}
\usepackage{enumitem}
\usepackage{yhmath}
\usepackage{derivative}
\usepackage{physics}

\usepackage{selectp}

\begin{document}

\section*{Dimostrazioni di Conservativit\`a e Calcolo dei Fattori di Scala}

\subsection*{1. Caso Cilindrico (Simmetria Cilindrica)}
Le coordinate cilindriche $(r, \theta, z)$ sono legate alle coordinate cartesiane $(x, y, z)$ dalle relazioni:
\begin{equation*}
    x = r\cos\theta, \quad y = r\sin\theta, \quad z = z
\end{equation*}
Il vettore posizione $\vec{r}$ si scrive quindi come:
\begin{equation*}
    \vec{r} = r\cos\theta\hat{i} + r\sin\theta\hat{j} + z\hat{k}
\end{equation*}

\subsubsection*{Calcolo dei Fattori di Scala}
I fattori di scala $h_i$ si calcolano come il modulo del vettore derivato rispetto a ciascuna coordinata:
\begin{itemize}
    \item \textbf{Fattore $h_r$:}
    \begin{equation*}
        \frac{\partial \vec{r}}{\partial r} = \cos\theta\hat{i} + \sin\theta\hat{j} \implies h_r = \left| \frac{\partial \vec{r}}{\partial r} \right| = \sqrt{\cos^2\theta + \sin^2\theta} = 1
    \end{equation*}
    \item \textbf{Fattore $h_\theta$:}
    \begin{equation*}
        \frac{\partial \vec{r}}{\partial \theta} = -r\sin\theta\hat{i} + r\cos\theta\hat{j} \implies h_\theta = \left| \frac{\partial \vec{r}}{\partial \theta} \right| = \sqrt{(-r\sin\theta)^2 + (r\cos\theta)^2} = r
    \end{equation*}
    \item \textbf{Fattore $h_z$:}
    \begin{equation*}
        \frac{\partial \vec{r}}{\partial z} = \hat{k} \implies h_z = \left| \frac{\partial \vec{r}}{\partial z} \right| = \sqrt{1} = 1
    \end{equation*}
\end{itemize}

\subsubsection*{Metodo A: Test del Rotore}
Sia una forza radiale a simmetria cilindrica $\vec{F} = F_r(r)\hat{u}_r$. Impostiamo la matrice col fattore esterno $1/(h_r h_\theta h_z) = 1/r$:
\begin{equation*}
    \vec{\nabla} \times \vec{F} = \frac{1}{r} 
    \begin{vmatrix} 
        \hat{u}_r & r\hat{u}_\theta & \hat{u}_z \\ 
        \frac{\partial}{\partial r} & \frac{\partial}{\partial \theta} & \frac{\partial}{\partial z} \\ 
        F_r(r) & 0 & 0 
    \end{vmatrix} = \vec{0}
\end{equation*}
Le derivate parziali di $F_r(r)$ rispetto a $\theta$ e $z$ sono nulle poich\'e la forza dipende solo da $r$.

\subsubsection*{Metodo B: Lavoro Infinitesimo}
Lo spostamento infinitesimo \`e $d\vec{r} = h_r dr\hat{u}_r + h_\theta d\theta\hat{u}_\theta + h_z dz\hat{u}_z = dr\hat{u}_r + rd\theta\hat{u}_\theta + dz\hat{u}_z$. Il lavoro infinitesimo risulta:
\begin{equation*}
    dW = \vec{F} \cdot d\vec{r} = \big(F_r(r)\hat{u}_r\big) \cdot \big(dr\hat{u}_r + rd\theta\hat{u}_\theta + dz\hat{u}_z\big) = F_r(r)dr
\end{equation*}
L'integrale dipende solo dalle posizioni radiali iniziale e finale, provando la conservativit\`a.

\newpage

\subsection*{2. Caso Sferico (Forze Centrali)}
Le coordinate sferiche $(r, \theta, \phi)$ sono legate alle cartesiane da:
\begin{equation*}
    x = r\sin\theta\cos\phi, \quad y = r\sin\theta\sin\phi, \quad z = r\cos\theta
\end{equation*}
Il vettore posizione \`e $\vec{r} = r\sin\theta\cos\phi\hat{i} + r\sin\theta\sin\phi\hat{j} + r\cos\theta\hat{k}$.

\subsubsection*{Calcolo dei Fattori di Scala}
\begin{itemize}
    \item \textbf{Fattore $h_r$:}
    \begin{equation*}
        \frac{\partial \vec{r}}{\partial r} = \sin\theta\cos\phi\hat{i} + \sin\theta\sin\phi\hat{j} + \cos\theta\hat{k} \implies h_r = \sqrt{\sin^2\theta(\cos^2\phi+\sin^2\phi)+\cos^2\theta} = 1
    \end{equation*}
    \item \textbf{Fattore $h_\theta$:}
    \begin{equation*}
        \frac{\partial \vec{r}}{\partial \theta} = r\cos\theta\cos\phi\hat{i} + r\cos\theta\sin\phi\hat{j} - r\sin\theta\hat{k} \implies h_\theta = \sqrt{r^2\cos^2\theta(\cos^2\phi+\sin^2\phi)+r^2\sin^2\theta} = r
    \end{equation*}
    \item \textbf{Fattore $h_\phi$:}
    \begin{equation*}
        \frac{\partial \vec{r}}{\partial \phi} = -r\sin\theta\sin\phi\hat{i} + r\sin\theta\cos\phi\hat{j} + 0\hat{k} \implies h_\phi = \sqrt{r^2\sin^2\theta(\sin^2\phi+\cos^2\phi)} = r\sin\theta
    \end{equation*}
\end{itemize}

\subsubsection*{Metodo A: Test del Rotore}
Sia la forza centrale $\vec{F} = F_r(r)\hat{u}_r$. Inseriamo lo Jacobiano esterno $1/(h_r h_\theta h_\phi) = 1/(r^2\sin\theta)$:
\begin{equation*}
    \vec{\nabla} \times \vec{F} = \frac{1}{r^2 \sin\theta} 
    \begin{vmatrix} 
        \hat{u}_r & r\hat{u}_\theta & r\sin\theta\hat{u}_\phi \\ 
        \frac{\partial}{\partial r} & \frac{\partial}{\partial \theta} & \frac{\partial}{\partial \phi} \\ 
        F_r(r) & 0 & 0 
    \end{vmatrix} = \vec{0}
\end{equation*}
Essendo $F_r$ indipendente dagli angoli, le derivate parziali rispetto a $\theta$ e $\phi$ svaniscono.

\subsubsection*{Metodo B: Lavoro Infinitesimo}
Lo spostamento infinitesimo \`e $d\vec{r} = dr\hat{u}_r + rd\theta\hat{u}_\theta + r\sin\theta d\phi\hat{u}_\phi$. Il prodotto scalare fornisce:
\begin{equation*}
    dW = \vec{F} \cdot d\vec{r} = F_r(r)dr \implies W_{A \to B} = \int_{r_A}^{r_B} F_r(r)dr
\end{equation*}
L'indipendenza del lavoro dalle coordinate angolari $\theta$ e $\phi$ chiude la dimostrazione.

\end{document}